\mxcorrections{2}\label{C02-CORR}
\textbf{Usage professeur.} Les réponses qui suivent concernent les tâches non résolues dans le cours. Faire d’abord expliciter la démarche ; fournir une aide graduée avant la solution. Les durées sont des estimations à étalonner en classe.
\subsection*{Diagnostic et remédiation}
\begin{enumerate}
\item $x^2-6x+9$. \textit{Action :} Reconstituer le produit de deux binômes.
\item $(x-4)(x+4)$. \textit{Action :} Utiliser la différence de deux carrés sur deux exemples.
\item $x=2$ ou $x=-5$. \textit{Action :} Faire verbaliser la règle du produit nul et son hypothèse.
\item $x=-3$ ou $x=3$. \textit{Action :} Distinguer la racine carrée positive des solutions de l’équation.
\item Négatif avant 3, nul en 3, positif après 3. \textit{Action :} Tracer la droite puis faire le tableau de signe.
\item 4 atteint uniquement en 1. \textit{Action :} Partir du caractère positif d’un carré.
\end{enumerate}
\subsection*{Activité de recherche}
Les longueurs positives x et y vérifient $x+y=10$ et $xy=24$. Elles sont les racines de $t^2-10t+24=(t-4)(t-6)$. Dimensions 4 m et 6 m ; contrôle de l’aire et du périmètre.
\subsection*{Applications immédiates}
\begin{enumerate}
\item $g(x)=(x+1)^2-4=(x-1)(x+3)$. Racines 1 et $-3$, minimum $-4$ en $-1$.
\item La première est $3(x-1)^2=0$, donc 1. Pour la seconde, $\Delta=-16<0$ : aucune solution réelle.
\item $f(x)=a(x-1)(x-4)$ ; $4a=-8$, donc $f(x)=-2(x-1)(x-4)$.
\item $(x-5)(x+1)<0$ donne $\intoo{-1}{5}$. Le sommet de $x(8-x)$ est en 4, hors du domaine ; sur $\intff{0}{3}$, la fonction croît et son maximum vaut 15 en 3.
\end{enumerate}
\subsection*{Exercices : démarche et résultats}
\paragraph{Exercice 2.1 -- Compléter un carré.} $x^2-8x+7=(x-4)^2-9=(x-1)(x-7)$. Le coefficient dominant est positif ; le trinôme est négatif ou nul entre les racines : $\intff{1}{7}$.
\paragraph{Exercice 2.2 -- Discuter sans calcul lourd.} $x^2+4>0$ toujours, sans racine. $(x+2)^2\geq0$, racine double $-2$. Pour le dernier, racines $-4$ et 1 ; positif sur $\intoo{-4}{1}$, négatif à l’extérieur, nul aux bornes.
\paragraph{Exercice 2.3 -- Choisir une stratégie.} Identité : $x=\pm3$. Somme et produit : 4 et 5. Pour la troisième, pas de racine entière évidente : $\Delta=17$ et $x=\dfrac{3\pm\sqrt{17}}{4}$. Le discriminant convient dans tous les cas, mais n’est pas toujours le plus efficace.
\paragraph{Exercice 2.4 -- Retrouver une parabole.} $f(x)=a(x+2)(x-3)$. En 1 : $-6a=12$, donc $a=-2$ ; $f(x)=-2x^2+2x+12$, d’où $f(0)=12$.
\paragraph{Exercice 2.5 -- Un potager contre un mur.} Les deux largeurs utilisent $2x$ mètres ; longueur $20-2x$. Domaine $\intff{0}{10}$ pour la fonction, $\intoo{0}{10}$ pour un rectangle non dégénéré. $A(x)=20x-2x^2=50-2(x-5)^2$. Maximum 50 m$^2$, dimensions 5 m et 10 m.
\paragraph{Exercice 2.6 -- Un paramètre.} $\Delta=4-4m$. Deux racines si $m<1$ ; racine double 1 si $m=1$ ; aucune si $m>1$. Le paramètre m reste fixé quand on résout en x.
\paragraph{Exercice 2.7 -- Lien avec les suites.} $u_n=(n-1)(n-4)$ ; négatif pour $1<n<4$. Comme n est entier, les seuls rangs sont 2 et 3. L’ensemble des rangs n’est pas l’intervalle réel entier.
\paragraph{Exercice 2.8 -- Reprise différée.} Le minimum vaut $-5$ en $-3$. L’inéquation donne $(x+3)^2\leq1$, soit $-1\leq x+3\leq1$, donc $\intff{-4}{-2}$.
\paragraph{CQFD, exercice 43 p. 99, questions a et c.} a : racine évidente 1, produit $-30$, donc seconde racine $-30$. c : écrire $4x^2+9x+5$ ; racine $-1$, produit $\dfrac54$, donc seconde racine $-\dfrac54$. Vérifier les sommes, respectivement $-29$ et $-\dfrac94$.
\paragraph{CQFD, exercice 49 p. 99, questions a et b.} Les nombres sont les racines de $t^2-st+p$. a : $t^2-t+1$, discriminant $-3$ : aucun couple réel. b : $t^2+60t+756$, discriminant $576=24^2$ : $-18$ et $-42$. Leur somme vaut $-60$ et leur produit 756.
\subsection*{QCM : réponses et raisons}
\textbf{QCM 1.}
\begin{enumerate}
\item C. Annuler chaque facteur.
\item A. $4-12=-8$.
\item D. Le carré s’annule en 2, valeur 5.
\item B. $-\dfrac{b}{a}=\dfrac{6}{2}=3$.
\item C. Coefficient dominant positif, inégalité stricte.
\end{enumerate}
\textbf{QCM 2.}
\begin{enumerate}
\item D. Identité remarquable.
\item A. $\dfrac{c}{a}=-\dfrac{2}{3}$.
\item C. La fonction croît jusqu’à 2 ; sur ce domaine, bord droit 1.
\item B. Pour a non nul, le discriminant nul caractérise la racine double.
\item D. Le coefficient multiplicateur est libre et non nul.
\end{enumerate}
\subsection*{Vrai ou faux}
\begin{enumerate}
\item Faux : $x^2+1$ est strictement positif.
\item Faux sans hypothèse sur a ; $-(x-1)(x-2)$ est positif entre 1 et 2.
\item Vrai : le carré d’un réel est positif ou nul ; égalité seulement en 2.
\end{enumerate}
\paragraph{Différenciation.} Pour un blocage technique, proposer la première étape seulement (domaine, facteur ou univers de référence), puis retirer l’aide. Pour un élève autonome, demander une seconde méthode et le traitement d’un cas limite. À l’oral, autoriser une reformulation initiale dans la langue maîtrisée, puis exiger une phrase mathématique précise en français.
